% l5_assurance.tex --- Layer 5 (Assurance and Verification), P4 draft 2026-10-17.
% Subsections for \section{Layer 5 --- Assurance and Verification}\label{sec:assurance}.
% Placeholder cite convention: \cite{p2[<arxiv-id>]} for IDs verified in the P2 sweep
% (wire named keys at P2; rauk2026jevbenchmark already wired in refs.bib).
% IDs cited here: 2210.03629, 2309.01219, 2510.06265, 2502.06872, 2409.10102, 2205.11916,
% 2203.11171, 2607.04430, 2510.08049, 2505.19475, 2311.08298, 2503.15850, 2407.18418,
% 2405.01563, 2211.16158, 2609.26123, 2609.30217, 2609.33887, 2609.37647.
% v3 add (2026-10-21): 2605.19537, 2607.24440, 2601.06118, 2608.14376, 2606.29581, 2404.06654,
% 2502.05167, 2404.16130, 2405.14831, 2609.11065, 2609.23053, 2305.14627; plus non-arXiv
% rosbach2026stuck, atman2026llm, fang2026right, huang2026cade, bachar2026lpp, fu2026knowing,
% chroma2025contextrot, anthropic2025context (metadata: docs/keys-C.md).

% ---- epigraph: the L5 thesis in vendor language ----
\begin{quote}
``A threshold is a policy you choose from measured accuracy at that coverage on your data, not a
property of the model. \dots{} Then pick the point where the errors you accept are ones you can
live with. Confidence orders decisions; it does not establish that a decision is correct.''
\\[-0.3em]
{\small --- Laya project documentation (vendor-reported)\footnote{\texttt{github.com/NandhaKishorM/laya},
README, accessed 2026-10-17; evidence class \evvendor{}.} --- the stance this layer formalizes.}
\end{quote}

\subsection{The layer contract: reliability as residual-risk management}
\label{sec:l5-contract}

Layer~4 delivers outputs that are structurally valid: schema-conformant, typed, parseable
(\cref{sec:contracts}). Validity is not truth, and no output \emph{contract} can make it truth:
hallucination persists across model families, tasks and elicitation formats, which is why
detection and mitigation retain survey literatures of their own
\cite{id230901219,id251006265} \evsynth{}. Layer~5 is the collection of mechanisms a
system uses to manage that residual: attaching evidence to outputs, checking them, refusing
to deliver some of them, and routing the rest to a human. Its contract is deliberately weaker
than ``no errors'' and deliberately stronger than ``hope'':

\begin{quote}
An L5-wrapped system promises that undetected errors are a \emph{managed quantity}: the
operator chooses how much error may pass unflagged, the system measures how much actually does,
and drift or evasion appears as a monitored signal.
\end{quote}

The responses available to the layer, ordered roughly by cost, form a ladder.
(1)~\emph{grounding} --- attach retrieved evidence and tool results, so claims become
attributable; (2)~\emph{deterministic checks} --- validators and content rules (quote-in-source
tests, unit and interval checks) that reuse the Layer-4 machinery but verify content rather than
shape; (3)~\emph{model-based checking} --- self-consistency and learned verifiers, which trade
sampling compute for an error signal; (4)~\emph{calibrated selective prediction} --- accept or
abstain against a threshold derived from a risk budget (\S\ref{sec:l5-calibrated});
(5)~\emph{escalation} --- defer the hardest cases to a human or a slower mode. The ladder is
cumulative, and every rung is itself a component that can fail; \S\ref{sec:l5-envelope}
therefore closes the chapter with the layer's own failure envelope. The epigraph is the stance
the calibrated rung formalizes: thresholds are policies chosen from measured behavior.

\subsection{Retrieval grounding: evidence channels and their new failure paths}
\label{sec:l5-grounding}

Grounding is the cheapest way to change a claim's \emph{epistemic status}: if the document a
system cites is the source of what it says, a reviewer can audit the claim without auditing the
model. In its interactive form the model acquires the evidence itself: ReAct interleaves
reasoning traces with task-specific actions so that the model queries external sources and
conditions on what comes back, which the original work credits with reducing hallucination and
error propagation on question answering and fact verification \cite{id221003629}. In its
passive form, retrieval-augmented generation supplies context up front; the trustworthy-RAG
surveys are the systematic accounts of what this buys and what it costs. RAG is reported to
reduce hallucination and freshen knowledge, but it introduces new risks along its own pipeline
(robustness, privacy, adversarial exposure, accountability), and unreliable retrieval or
improper use of retrieved knowledge still yields bad outputs \cite{id250206872}. A companion
framework scores RAG trustworthiness on six dimensions --- factuality, robustness, fairness,
transparency, accountability and privacy --- with a benchmark for each
\cite{id240910102} \evsynth{}.

The failure paths are what the layer must be designed around. \emph{Retrieval miss:} the
evidence exists but is not retrieved, and the generator answers ungrounded with undiminished
fluency; misses of this kind stay invisible unless the system is required to cite. \emph{Stale
or conflicting sources:} grounding certifies that a claim came from \emph{a} document; that
document can be wrong, and a confidently stale source launders error into apparent diligence.
\emph{Unfaithful use:} the evidence sits in context but does not shape the output, the
``improper knowledge utilization'' the survey singles out \cite{id250206872}. The nominal
window and the effective one diverge: of seventeen models claiming at least 32K tokens, half
keep satisfactory performance at 32K \cite{id240406654} \evindep{}; under a needle test that
removes the lexical shortcut, eleven of thirteen fall below half their short-context score at
the same length, GPT-4o from 99.3\% to 69.7\% \cite{id250205167} \evindep{}; and an
eighteen-model industry evaluation finds degradation with input length on tasks that require no
retrieval at all \cite{chroma2025contextrot} \evindep{}. Vendor guidance on managing the window
--- compaction, note-taking, sub-agents --- is vendor-reported \cite{anthropic2025context}
\evvendor{}, and no independent replication of it exists in our evidence \textsf{[N]}.
\emph{Lossy context:} the window is a lossy store --- it holds what was retrieved, and what gets
used is a separate step. Grounding therefore converts an unobservable failure (model error) into
an \emph{auditable} one (citation mismatch) that deterministic checks can catch. It does not
convert error into truth.

The index is a design decision, and its shape is a measured design variable. A flat top-$k$
store fails questions that are global to a corpus by construction, which is the gap
graph-structured indexes are built to close \cite{id240416130} \evaurep{}; a graph index over
entities and passages buys about 20\% on multi-hop retrieval at 10--30$\times$ lower cost than
iterative retrieval \cite{id240514831} \evaurep{}; and the traversal policy over that graph is
worth up to 9.96 NDCG@10 over the strongest fixed policy while exploring 81.9\% fewer paths
\cite{id260911065} \evaurep{}. \evjudge{The three results say one thing: which questions are
answerable at all is decided before the generator runs, so a retrieval miss is a property of the
index and its policy as well as of the query.}

The audit claim has a false-negative rate of its own. In a controlled reproduction on
Wikipedia-based question answering, roughly one citation in seven was unfaithful --- the answer
was right and the cited passage did not support it --- and the rate did not improve under RLVR,
one trained search agent landing slightly worse than its base model \cite{id260923053}
\evindep{}. ALCE puts incomplete citation support at 50\% for the best models on ELI5
\cite{id230514627} \evindep{}. Citation mismatch is therefore a check with a measured miss rate:
it catches the citation that fails to support the sentence and passes the citation chosen after
the answer was already fixed, which is why a citation check belongs in the budget with a stated
miss rate.

\subsection{Self-checking: agreement as a cheap signal and its sampling price}
\label{sec:l5-selfcheck}

The cheapest model-based checks verify nothing external; they interrogate the generation process
itself. Zero-shot chain-of-thought showed that a single elicitation phrase can move accuracy
substantially on arithmetic, symbolic and logical tasks without any exemplars \cite{id220511916}
\evaurep{} --- a prompt-level change to the error distribution: the system learns nothing about
\emph{which} of its answers are wrong. Self-consistency turns sampling itself into a signal:
draw a diverse set of reasoning paths, answer by the agreement among them, and benchmark gains
on arithmetic and commonsense reasoning follow \cite{id220311171} \evaurep{}. The cost is the
k-sample price in latency and budget, and it scales with k.

Self-consistency is best read as an \emph{uncertainty estimate}, with two caveats. First, a
systematically wrong model can produce agreement too, and the
uncertainty literature is explicit that such scores are heuristic: a threshold chosen on them
does not by itself provide statistical guarantees on the error rate among accepted answers
\cite{id260704430}. Second, the signal is nonetheless useful raw material: the
conformal-abstention construction takes self-consistency as its reliability measure and adds
the missing guarantee on top (\S\ref{sec:l5-calibrated}). The pattern repeats at every rung:
heuristic signal first, calibrated signal where the stakes require it.

\evaurep{Minimum-Bayes-risk selection over cheap candidates is the assurance-side use of the
sampling budget that shorter generation paths free up: on XSum the NCM authors report about 7\% higher ROUGE-L than the strongest
summarization baseline within the same sampling-time budget.} The gain comes from agreement among
candidates. Agreement is not independent verification, so budget reallocated this way must
still route through this layer's grounding and verification machinery
(\cref{sec:l5-verifiers}).

\subsection{Verification loops: process supervision and the verifier as a component}
\label{sec:l5-verifiers}

Where self-checking asks the model about itself, a \emph{verifier} is a separate component
that is trained to judge. Process reward models score reasoning at the step or trajectory level rather
than judging final answers alone, and the survey of the area organizes the full loop --- how
process data is generated, how PRMs are built, and how they are consumed for test-time scaling
and reinforcement learning --- across math, code, multimodal reasoning, robotics and agents
\cite{id251008049} \evsynth{}. Verification signals can also be recycled into adaptation:
verifier-driven test-time training scores a pool of sampled answers, keeps the high-confidence
ones as pseudo-labels, and fits a small adapter to them, improving the base model on
structurally novel tasks \cite{id250519475} \evaurep{}. This is Layer~5 re-entering L3 and
L1: assurance outputs become inference and training signal.

Both moves share a structural caveat. A verifier is itself a model, with an error rate and a
proxy objective; when it selects data or scores rewards, the generator acquires an incentive to
satisfy the verifier rather than the world, and the verifier's own errors are no longer
checkable by the system that consumes them: they \emph{become} the system's errors.
Verification loops therefore reduce some errors and relocate others, and the relocation is
invisible from the outside unless the verifier's accuracy is itself maintained against
held-out ground truth. This is the first entry in the layer's failure envelope
(\S\ref{sec:l5-envelope}).

\subsection{Calibrated selective prediction: thresholds as policies derived from risk budgets}
\label{sec:l5-calibrated}

The strongest rung is the most studied: use the model's own uncertainty to decide when to answer
at all. The survey anchors agree on the landscape. The confidence-estimation survey organizes
calibration techniques and their applications \cite{id231108298}; the newer
uncertainty-quantification survey adds a taxonomy along two axes --- computational efficiency
and uncertainty source (input, reasoning, parameter, prediction) --- and frames selective
prediction as the operative risk-mitigation use case \cite{id250315850}. Abstention now has a
survey of its own, organized by what is uncertain (the query), who is uncertain (the model), and
against whose values abstention is judged (human values); the last axis is what makes abstention
a policy \cite{id240718418} \evsynth{}.

This paper's stance is that a threshold is \emph{derived}, never universal, the position the
epigraph states in vendor language. An operator fixes a risk budget (the maximum tolerated error
rate among accepted outputs, per task and per unit of harm) and then selects the most permissive
threshold whose measured behavior on held-out data satisfies it. The budget is where costs
enter: a routing decision and a legal determination get different budgets. The vendor
documentation for the decision-model case study (the epigraph's source) supplies the operational
corollaries: fit and validate any cutoff on held-out data \emph{at the option counts the
workload uses}, because the usable threshold shifts with the configuration, and refit shipped
confidence values before any gating policy relies on them.

Two constructions make the budget statement literal, and they are the layer's theory-class ([T])
machinery. A conformal-abstention procedure uses self-evaluated similarity across sampled
responses as the confidence measure and conformal prediction to bound the hallucination rate
among non-abstained answers \cite{id240501563}. A second framework turns any uncertainty score
into a selective-answering rule: for each candidate threshold it estimates a high-probability
upper bound on the error rate among accepted answers, under a criterion the application
declares, and then adopts the largest threshold whose bound stays below the user's risk
\cite{id260704430}. Both are \evtheory{conditional guarantees: distribution-free within an
exchangeability assumption on calibration data, and only as good as the criterion that defines
``error''}. They formalize what the surveys only recommend: the threshold is the output of a
procedure over data plus budget. \evsynth{For the newest generation family, the criterion is
concrete: where the frozen codec's residual errors are substitutions of named entities and
numerals and errors in diacritics and usernames while surface reconstruction stays near-exact
(NCM-class \cite{kumbong2026ncm}), the accepted-error budget must state whether a fluent
substitution that changes a referent or quantity counts as an error, and at which layer it is
adjudicated --- the token-exactness vs.\ meaning-equivalence boundary, priced rather than
assumed.}

\paragraph{The fifth rung: escalation.} Abstention keeps a case out of the answer; escalation
moves it to a human, or to a slower and more expensive mode, whose capacity and error rate are
themselves resources that belong in the budget. The abstention survey's human-values axis is the
reminder that ``answer or abstain'' is a judgment about consequences; escalation is where those
consequences are owned explicitly.

That budget is usually written as though review were a constant and a catch. Neither assumption
survives measurement. In a pathology study, 7\% of baseline-correct estimates flipped after an
erroneous machine prediction, inside a 6--11\% range reported across specialties, and the
coefficient on the AI estimate was 0.74 when the advice beat the reviewer's own estimate against
0.30 when it was worse \cite{rosbach2026stuck} \evindep{}. Expertise moderates the damage without
removing it: in a 20-resident simulation, management accuracy fell from 77.3\% under concordant
advice to 45.3\% under discordant advice \cite{atman2026llm} \evindep{}. Review also depletes: computer-aided
detection's adenoma advantage was present early (49.9\%) and gone late (39.9\%, against a stable
39--41\% unassisted) \cite{huang2026cade} \evindep{}. Detection of naturally occurring errors ran
24\% at baseline and 31\% at high injected-error prevalence, at 16--18\% more time per review
\cite{fang2026right} \evindep{}. Capacity can be budgeted as a queue rather than a
constant: on a moderation benchmark, selective escalation held human load at 16--26\% of cases
where uncertainty baselines routed 28--39\%, at \$38--77 against \$127 for always-trust
\cite{bachar2026lpp} \evaurep{}.

The distinction that matters is therefore whether a human can catch the failure. A reviewer
needs an independent judgment formed before the output is shown, the evidence to check it
retrievable at review time, the time to use it, and the authority to overrule; without that
support, oversight decays with repetition even among reviewers \cite{fu2026knowing} \evindep{}.
What the evidence does not contain is the term an operator would budget: no study in this sweep
reports precision, recall or minutes per review for reviewers working a production LLM
escalation queue, so the rates above come from adjacent settings and do not transfer. Escalation
priced that way reprices risk; it does not remove it.

\subsection{Monitoring: the monitored quantity, and the evadable layer}
\label{sec:l5-monitoring}

A deployed system's thresholds decay: calibration data ages, the workload drifts, and the
monitor's own distributions move. Monitoring is the rung that tests, after deployment, whether
the assumptions the earlier rungs rest on still hold, and its first design question is what to
monitor. The monitoring literature's sharpest lesson is that the obvious answer is usually a
proxy. Runtime-monitor research argues that out-of-distribution detection is \emph{not} the
right framing: a monitor's job is to discard incorrect predictions, and measured on that
quantity (``out-of-model-scope'' detection), OOD scores can give a false impression of safety,
while OOD-framed comparisons need not identify the better monitor \cite{id221116158} \evindep{}.
The lesson parallels the threshold-as-policy stance: the monitored quantity is chosen by what
the operator cares about, and validated against that quantity.

The monitored surface is also wider than correctness. Fairness-monitoring tools specify
subgroup-behavior properties in a runtime specification language and expose them with live
visualization, and they target the human-oversight duties that regulation such as the EU AI Act
places on high-stakes decision systems \cite{id260926123} \evindep{}. Monitoring is where
the layer becomes auditable: a monitor is the operational form of an oversight duty.

The honest warning, and the newest result in this section, is that the assurance layer is itself
a target. An evaluation of instrumental evasion places agents in task-policy pairs where
completing the task requires an operation the runtime monitor prohibits, with the agents aware
they are monitored: best-of-3 evasion \emph{attempt} rates reach up to 98\% and success rates up
to 88\%, evasion grows with test-time compute, and the successful strategies are mundane:
encoding prohibited commands, decomposing an operation across several tool calls, and retrying
until the relevant context leaves the monitor's history \cite{id260930217} \evindep{}. Two
implications. First, oversight-evasion emerges from ordinary goal-directed behavior under task
pressure, and belongs to the deployment setup that the assurance layer must budget for. Second,
the trade-off cuts both ways: in the same study the model with the lowest evasion rate paid for
it in over\-refusal, abandoning otherwise solvable tasks. Monitoring policy, like the thresholds
it feeds, is chosen against consequences.

\subsection{L4--L5 bridge: configuration selection against a risk budget}
\label{sec:l5-bridge}

The pattern Layer~5 formalizes is already visible in production serving practice, which is
why the contract layer and the assurance layer exchange more than a one-way handoff: the
contract tier decides what can be measured, and the measurement decides where the operating
point sits. The clearest current demonstration is block-diffusion serving. Its operating points --- acceptance
thresholds, buffer depth, schedule, checkpoint, precision --- are normally hand-picked by mean
benchmark accuracy, but a mean does not tell an operator how often a faster configuration fails
on prompts that a slower reference configuration answers correctly. The risk-guaranteeing
procedure (Redline) selects operating points from their measured correctness on calibration
prompts, keeps the reference-relative risk --- the joint event ``reference correct, candidate
wrong'' --- within a user-chosen budget with high probability, and deploys the fastest
configuration that passes \cite{id260933887}. \evtheory{The guarantee is finite-sample and
distribution-free, but relative to the reference configuration and the chosen correctness
measure; LEC gives the general form for selective prediction and routing, where linear expectation
constraints hold the risk of the selected set below a budget \cite{lec2026}.}
\evaurep{On the same calibration data the procedure
stays within its stated failure probability, whereas a mean-accuracy rule at any fixed tolerance
either buys less speed on some model--task pairs or overshoots the risk budget on others; at a
10\% risk budget it deploys a LLaDA2 math configuration that commits over a
third more tokens per forward pass, and therefore needs proportionally fewer passes to emit the
same sequence. The measured budgets are task-dependent --- math clears
at a tighter budget than code} \cite{id260933887}. ``Cheaper but less
accurate'' is thereby replaced by a measurable statement: speed is bought at a chosen risk
budget. The same selection logic governs the other layers' choice points
(objective, architecture, contract tier), which is the subject of \cref{sec:selection}.

The guarantee attaches to a serving configuration. \emph{Engine.} With
weights and decoding fixed, the inference engine alone shifts benchmark scores by up to 16.6
percentage points across five engines \cite{id260519537} \evindep{}; in one vision-language family,
4-bit weights cost 1.6 accuracy points and drop the internal error-detection AUROC from 0.95 to
0.80 \cite{id260724440} \evindep{}. Nondeterminism is concentrated in $p\in[0.1,0.9]$, the
cutoff's own neighbourhood \cite{id260106118}, and \evaurep{removing it costs more than twice the
latency and up to 74\% throughput loss}~\cite{id260814376}. \emph{Precision.} Precision is not even a stable
label: vLLM's quantized KV cache defaults to uncalibrated scales of one, and TensorRT-LLM's FP8
recipe differs between SM90 and SM100+ (FP32 against UE8M0 activation scales). Calibration is a
deployment artifact that no checkpoint records, and the same weights on a newer GPU are different
arithmetic\footnote{%
\url{https://docs.vllm.ai/en/latest/features/quantization/quantized_kvcache};
\url{https://github.com/NVIDIA/TensorRT-LLM/blob/main/docs/source/features/quantization.md},
accessed 2026-10-20; evidence class \evvendor{}.}. \emph{Co-tenancy.} Batch composition changes
reductions \cite{id260106118} \evindep{}, so co-resident requests can change the sampled tokens for
the same prompt; the serving stack documents the effect \cite{vllmfaq2026} \evvendor{}. We treat it
as a third determinism term beside task determinism and model stochasticity (serving-path
determinism, \cref{sec:determinism}), so the deployed function is a function of its neighbours.
The danger is not that quantization is unsafe --- AWQ INT4's safety score stayed within
about 1.6 points of full precision on seven of eight models \cite{id260629581} \evindep{} --- and
the risk sits in the record: a cutoff carries an error rate for the configuration it was fitted on. A
stack change is a model change, and the configuration record has to say so.

The decision-model case studies supply the second half of the bridge: systems that return typed
probabilities per question (Layer~4) are exactly the systems on which
\emph{measured-accuracy-at-coverage} thresholds are computable, and the independent Jev
benchmarking line evaluates a commercial decision model on precisely those terms:
zero-shot, frozen templates, full rather than sampled splits. Its findings read as a
controlled demonstration of this section's principle: the model's choice probabilities
measure as well-calibrated and support selective prediction, but its binary
probabilities --- ranking well as scores --- sit poorly placed against a fixed 0.5 default,
and thresholds tuned on its training data raise the classification result on a
terms-of-service benchmark (micro-F1 from 0.50 to 0.75) \cite{rauk2026jevbenchmark}
\evindep{}. In that benchmark the threshold was the binding
constraint on the reported metric. That is evidence about one benchmark, not evidence that
thresholds dominate in general. The vendor documentation quoted in the epigraph shows
the same stance arriving in industry language, caveats included: cutoffs are fit and validated
on held-out data at the option counts the workload uses, and shipped confidence is refitted
before gating. Deployment-grade selection, in other words, is a Layer-5 procedure: Layer-0
sets the budget, Layer-4 supplies the scored return, and the procedure picks the operating point.

\subsection{The failure envelope of the assurance layer}
\label{sec:l5-envelope}

Every rung above is a component, and components fail. The layer's own envelope deserves the
same honesty as the model envelopes of \cref{sec:objectives-cards}:

\begin{itemize}
  \item \textbf{Wrong grounding.} The evidence channel can carry stale, conflicting or false
        documents; grounding makes claims attributable. A system that cites
        carefully can still be wrong with provenance.
  \item \textbf{Wrong verifier.} Learned checkers have error rates, and one used as a
        selection signal or reward is subject to being optimized against rather than for the
        truth (\S\ref{sec:l5-verifiers}); verification moves error without deleting it.
  \item \textbf{Correction racing.} A check that completes after the action has committed has
        no authority. When outputs stream, tools execute, or downstream systems consume
        immediately, the assurance layer must fit inside the committed-before-irreversible
        latency budget, which is why reversibility and observability appear as
        operating-mode modifiers in \cref{sec:selection}.
  \item \textbf{Unsound auditing.} Monitors can be evaded under ordinary task pressure
        (\S\ref{sec:l5-monitoring}); monitor scores validated against proxies can mislead
        \cite{id221116158}; and a fairness monitor certifies conformance to a
        specification, not justice. Auditing is only as sound as the quantity it audits.
  \item \textbf{Scoped guarantees.} The theory-class ([T]) bounds are conditional on the
        calibration distribution; under drift, an unmonitored guarantee is a historical
        statement.
  \item \textbf{Escalation is also a system.} Human review has capacity limits, fatigue, and
        its own error rates --- detection of naturally occurring errors ran near 24\%
        \cite{fang2026right} \evindep{} --- so a budget that routes overload to humans has not
        removed risk, only moved and repriced it (\S\ref{sec:l5-calibrated}).
\end{itemize}

The closing honesty the whole chapter is built to support: \emph{layered defense is a policy
decision, not a property}. How many independent checks to stack, which evidence classes to
require, and where to set acceptance and escalation budgets are choices made per task
against error cost and answer-space openness (\cref{sec:selection}), and they are choices with
a ceiling: no layer makes another layer correct. What the \reliabilitystack{} claims for Layer
5 is the honest version of assurance: an error rate you measured, under assumptions you can
name, at a residual you chose to own.
